\documentclass[12pt]{article}
\usepackage{graphicx} %package to manage images
\usepackage{accsupp}     % for actual PDF alt text
\usepackage{accessibility}
\graphicspath{ {./figures/} }
\usepackage{caption}
\usepackage[font=scriptsize]{subcaption}
\captionsetup[figure]{labelsep=none}
\captionsetup[table]{labelsep=none}
\usepackage{bbm}
\usepackage{amsmath}
\usepackage{import}
\usepackage{array}
\usepackage{booktabs}     % no extra vertical space between paragraphs
\usepackage{indentfirst}
\usepackage{afterpage}
\usepackage{floatrow}
\usepackage{pdflscape}
\usepackage{soul}
\usepackage{float}
\usepackage{adjustbox}
\usepackage{longtable}
\usepackage{caption}
\usepackage{setspace}
\usepackage{afterpage}
\usepackage[margin=1in]{geometry}
\usepackage[round]{natbib}
\usepackage{hyperref}
\usepackage{titlesec}
\usepackage{threeparttable}
\usepackage{setspace}
\usepackage{float}
\floatstyle{plaintop}
\restylefloat{table}
\usepackage{rotating}
\usepackage{ragged2e}
\usepackage{authblk}



\usepackage{totcount}
\newtotcounter{citnum} % Define a new counter named 'citnum'

% Optional styling for nicer author/affil typography
\renewcommand\Authfont{\normalsize\bfseries}
\renewcommand\Affilfont{\small\normalfont}
\setlength{\affilsep}{0.3em} % space between author and affiliation

% \author[1,2,3]{Diane Coffey\thanks{Corresponding author: coffey@utexas.edu}}

% \author[1]{Siddhant Pandit}

% \affil[1]{Population Research Center, UT Austin}
% \affil[2]{Department of Sociology, UT Austin}
% \affil[3]{r.i.c.e., a research institute for compassionate economics}


\titleformat{\section}
  {\normalfont\bfseries\fontsize{14}{16}\selectfont}  % Bold 14pt
  {}{0pt}{}  % No numbering

\titleformat{\subsection}
  {\normalfont\bfseries\fontsize{12}{14}\selectfont}  % Bold 12pt
  {}{0pt}{}  % No numbering

\titleformat{\subsubsection}
  {\normalfont\itshape\fontsize{12}{14}\selectfont}  % Italic 12pt
  {}{0pt}{}  % No numbering
  
\begin{document}
\RaggedRight


\setlength{\parindent}{0.25cm}   % first-line paragraph indent
\setlength{\parskip}{0.5em}       % no extra vertical space between paragraphs

 \begin{center}
\Large{Disparities in Maternal Nutrition among Indian Social Groups: The Roles of Fertility Patterns and Household Wealth}      
 \end{center}
%\date{\today} 

%\maketitle
    

\singlespacing

\begin{abstract}
Maternal undernutrition is a threat to women’s health and a leading cause of child mortality and morbidity in India. This paper uses nationally representative data to estimate the fraction of women who begin pregnancy underweight -- a key indicator of maternal nutrition -- and presents the first estimates of the prevalence of prepregnancy underweight by caste and religion.  We use a nonparametric reweighting technique to show that more than one in five Indian women start pregnancy underweight.  Adivasis and Dalits are 55\%, and 40\% more likely, respectively, to be underweight before pregnancy than Forward Caste Hindus. Differences in parity progression and birth spacing do not explain the maternal nutrition gaps that we observe.  Instead, differences in household wealth explain a large fraction of the gaps. Notably, we find that a similar fraction of Muslim and Forward Caste Hindu women start pregnancy underweight.  Our results suggest the need for monitoring of and health interventions during the prepregnancy period.
\end{abstract}

{\bf{Keywords:}} pregnancy, nutrition, fertility, wealth, India

\newpage
\doublespacing

\section*{Introduction}

The intergenerational transmission of health is perhaps nowhere more evident than it is during pregnancy. Nutrition is an important pathway: undernourished mothers are more likely to have low birth weight and preterm babies, who face a higher risk of death than babies born to better nourished mothers \citep{aizer2014intergenerational}.  

India has the largest absolute number of undernourished mothers and research has documented a disproportionately high burden of maternal undernutrition relative to the country’s level of economic development \citep{bhutta2000has, coffey2015prepregnancy}. Prematurity and low birth weight, both linked to maternal undernutrition, constitute the leading causes of child mortality \citep{fadel2017changes}. Among children who survive infancy, those born to undernourished mothers are more likely to grow up to be malnourished and to suffer from chronic disease as adults \citep{underwood2002health, almond2011killing}.

Prepregnancy underweight, defined as having a body mass index (BMI) below 18.5 at conception, is a population-level indicator of maternal undernutrition \citep{who_nlis_malnutrition_women_2021}. This paper contributes the most recent estimates of prepregnancy underweight for India, and the first estimates of prepregnancy underweight stratified by Indian social group.  It then uses Kitagawa decomposition to answer the question: Can differences in fertility patterns or differences in household wealth explain a larger fraction of social group differences in prepregnancy underweight?

To estimate prepregnancy underweight from representative, cross-sectional data, we construct a nonparametric reweighting function to correct for selection into pregnancy.  This method was developed by \cite{coffey2015prepregnancy}, which found that, in 2005, 42\% of Indian women began pregnancy underweight.  Using nationally representative data from 2019-2021, we find that 20\% of Indian women begin pregnancy with a BMI below 18.5. Despite improvement, maternal nutrition remains an important population health challenge.  

We also contribute to a growing literature on how caste shapes population health  \citep{vyas2022social, gupta2020caste, coffey2019local, lopalo2019consequences, bora2019persistent}. We find that the burden of prepregnancy underweight is highest among Adivasi, or indigenous, women (27\%), followed by Dalits, a group of oppressed castes who have been treated as ``untouchable'' (23\%), and women from Other Backward Classes (20\%).  Muslim women who are not from a Dalit or Adivasi background have similar rates of prepregnancy underweight as Hindu women from dominant (``Forward'') castes (14\% and 16\%, respectively).  This is consistent with research that documents health advantages among Indian Muslims relative to Hindus, despite socioeconomic disadvantage \citep{guillot2010hindu, bhalotra2010puzzle, visaria2015preference, geruso2018neighborhood, sharma_gupta_2025}.  

In addition to documenting caste and religious differences, this study also investigates explanations. We first consider whether differences in parity progression and birth spacing can account for social group differences in prepregnancy underweight.  We test this hypothesis because prior research on maternal and infant health emphasizes the role of high parity and short birth spacing in depleting mothers' nutritional status \citep{winkvist1992new, bauserman2020relationship}.  Using a Kitagawa decomposition, we find that differences in parity and birth spacing between oppressed caste and Forward caste Hindu women account for between 4\% and 11\% of the gaps that we document.  

Second, we consider whether differences in household wealth can explain the gaps in prepregnancy underweight that we observe.  Using a Kitagawa decomposition, we find that the distribution of women into wealth quartiles explains about 68\% of the Adivasi-Forward caste Hindu gap in prepregnancy underweight, and about 62\% of the Dalit-Forward caste Hindu gap. Although differences in the distribution of women into wealth categories does explain a majority of the gaps we observe, meaningful fractions of the gaps remain unexplained.

In the next section, we present background information on maternal nutrition in India, on Indian social groups, and on candidate explanations for group-level variation in maternal nutrition.  We then describe our data source, the fifth round of the National Family Health Survey (NFHS) (2019-2021), India's Demographic and Health Survey (DHS).  We present the two complementary decomposition methods that we use to analyze maternal nutrition levels and gaps.  The first is a nonparametric reweighting technique that corrects for selection into pregnancy.  The second is a Kitagawa decomposition that helps us investigate candidate explanations for group-level differences in prepregnancy underweight.  Finally, we discuss our results and consider implications for policy and future research.


\section{Background}

\subsection{Maternal nutrition in India}

In the late 1990s, researchers noticed an ``Asian Engima:'' Asian children had higher rates of malnutrition than children in sub-Saharan Africa, despite being less poor, on average \citep{ramalingaswami1996commentary, underwood2002health}.  They posited that maternal malnutrition contributed to low birthweight and malnurition among children, and that maternal malnutrition was caused by low social status among Indian women.  

Subsequent research found empirical support for these theories.  Using data from 2005, \cite{coffey2015prepregnancy} found that 42\% of Indian women were underweight before pregnancy, compared to 17\% in sub-Saharan Africa.  Health guidelines recommend that underweight women gain 28–40 pounds during pregnancy to compensate for limited nutrient reserves \citep{yaktine2010weight, ismail2016gestational}. Yet, \cite{coffey2015prepregnancy} found that the average Indian woman gained so little weight during pregnancy that she weighed less at the end of pregnancy than the average woman in sub-Saharan Africa weighed at the beginning of pregnancy.  

Low gestational weight gain, which is common across low-income country settings \citep{wang2020gestational, perumal2023suboptimal}, makes the preconception period a critical window for nutrition interventions \citep{stephenson2018before}.  Yet, \cite{hathi2021women} find that many Indian women eat after the men in their households.   

\subsection{Indian social groups and population health}

Caste and religion are important axes of social inequality in India \citep{vaid2014caste, majid2015indian}. A person's \textit{jati}, or caste sub-group, influences their social interactions, marriage, occupation, and life chances \citep{deshpande2011grammar, thorat2012blocked}. Although caste is a concept closely associated with Hinduism, India's dominant religion, people from all religious backgrounds are stratified along caste lines. 

To address historical and ongoing discrimination, the central and state governments of India have affirmative action policies in higher education and government jobs. These policies are part of the Indian Constitution, which recognizes three historically disadvantaged groups: 

\begin{itemize}
\item Scheduled Castes (SC), comprising about 20\% of India's population, are \textit{jatis} that were traditionally considered ``untouchable.''  Both historically and today, Scheduled Castes do stigmatized work, such as work involving contact with dead bodies (leather work, cremation, removal of dead animals) or feces (cleaning latrines and bathrooms, maintaining sewers, emptying latrine pits).  People from Scheduled Castes refer to themselves as Dalit, or ``oppressed people.''  
\item Other Backward Classes (OBC) comprise about 43\% of India's population are are from \textit{jatis} considered socially, economically, or educationally deprived compared to the dominant, or ``Forward,'' castes.  ``Backward'' is a term used in the Constitution; we use the abbreviation OBC with no normative implications. 
\item Scheduled Tribes (ST), comprising 9\% of the total population, refer to India’s indigenous groups.  People from Scheduled Tribes face social and economic disadvantage due to geographic isolation and exclusion from state services.  To refer to Scheduled Tribes, we use the term Adivasi, a term of self-identification that means ``early inhabitants.''  
\end{itemize}

People from socially dominant \textit{jatis} are ineligible for affirmative action.  These \textit{jatis} are referred to as General castes or Forward castes.  We use the term Forward castes becasue it is used by the Indian Human Development Survey (IHDS), a large, nationally representative panel study of human development \citep{DesaiVannemanIHDS2011}.\footnote{Our main data source, the NFHS, codes General or Forward caste people as ``others.''}  

The fifth social group that we consider is Muslims, India's largest religious minority.  Muslims comprise about 16\% of the population.  Historically and today, Muslims face social exclusion and targeted violence.  Political movements linking Indian nationalism with Hinduism are central to the marginalization of Muslims \citep{farooqui2020political}. Muslims have ranked caste affiliations which range from extremely disadvantaged Dalit \textit{jatis} to socially dominant \textit{jatis} \citep{ahmad2007recognition}.  However, Muslims of all castes are typically ineligible for affirmative action \citep{hasan2018muslims, hasan2022wider}.   

The inequalities created by caste and religion shape population health in India. Dalits, Adivasis, and Muslims have lower life expectancies than Forward caste Hindus \citep{vyas2022social, gupta2022large}.  \cite{gupta2020caste} find more self-reported depression and anxiety among Dalits and Adivasis than among Forward caste Hindus.  \cite{coffey2019local} document gaps in child height among SC, ST, and OBC children when compared to children from Forward castes. Caste differences in health indicators are only partially explained by differences in socioeconomic status \citep{bora2019persistent, kumar2016explaining, vyas2022social, dutta2020effect}. 

However, social marginalization does not always correspond with health disadvantage.  Research finds better health for Muslim children and women on some indicators.  \cite{geruso2018neighborhood} show that Muslim infants are more likely to survive because they are exposed to better neighborhood sanitation, on average, than Hindu infants. \cite{sharma_gupta_2025} find that Muslim women are less likely to experience domestic violence than Hindu women, and that the difference can largely be explained by lower use of alcohol among Muslim men.  \cite{visaria2015preference} finds that selective abortion of girls is lower among Muslims than Hindus. 

\subsection{Fertility, birth spacing, and nutrition among Indian social groups}

Maternal health research about low- and middle-income country populations posits maternal depletion, the hypothesis that higher-order births and low spacing between births leave women nutritionally vulnerable \citep{winkvist1992new, king2003risk, shell2004maternal, bauserman2020relationship}.  Can differences in fertility and birth spacing explain differences in prepregnancy underweight across Indian social groups?  Whether differences in fertility and birth spacing will explain differences in prepregnancy underweight will depend on: 1) the magnitude of fertility and birth spacing differences across groups, and 2) the relationship between these variables and maternal nutrition in present-day India. 

Researchers have documented higher fertility among Muslims than Hindus \citep{dharmalingam2004pervasive, james2005accelerated}.  This difference has featured prominently in political discourse \citep{jeffery2002population, singh2020numbering}.  Each round of the NFHS finds higher fertility rates among Advisais, Dalits, and OBCs than among Forward castes, though caste differences in fertility are decreasing over time \citep{srinivas2024rapid}. The literature also notes differences in birth spacing across caste and religious groups, with Forward caste Hindu women being most likely to have more than two years between births \citep{kundu2025short, bhalotra2008birth}.  

However, \cite{coffey2021neonatal} provide suggestive evidence that, in India, higher parity may \emph{not} lead to maternal depletion.  They find that, controlling for siblingship size, higher-parity births are more likely to survive the first month of life than lower-parity births.  They posit that Indian women's social status, and consequently the balance of their food intake and work output, may be linked to age and parity progression, allowing women to gain weight across successive pregnancies. 

\subsection{Household wealth and nutrition among Indian social groups}

Both caste and religion influence a household's economic status.  \cite{tagade2020intergroup} estimated that, in 2013, Scheduled Castes, Scheduled Tribes, and Muslims each controlled 10\%, 6\%, and 9\% of India's total wealth, respectively.  These estimated wealth shares were considerably less than their 2013 population shares of 18\%, 9\%, and 12\%.  Other studies document social group differences in wealth and poverty as well \citep{thorat2017escaping}. 

There are several pathways through which household wealth may protect against undernutrition: it allows people to buy more and higher quality food \citep{raghunathan2021affordability}, they are more likely to have tools and technology that reduce their energy output \citep{deaton2009food}, and they may be able to afford cleaner water and better sanitation, which reduces energy loss to disease \citep{bhutta2008works}.  We therefore expect differences in household wealth to explain at least part of the gaps in prepregnancy underweight across social groups. 

\section{Data and Methods}
 
 In many high-income countries, pregnancy health monitoring systems allow researchers to track prepregnancy nutrition and study its relationship with maternal and birth outcomes \citep{cornish2024maternal}.  In India, and other low-income countries, however, researchers rely on large-sample, representative, cross-sectional surveys, such as the DHS, which measure women's heights and weights at the time of the survey.  

We use the most recent DHS from India, the National Family Health Survey (2019-2021) (NFHS-5). These data are publicly available from \texit{www.dhsprogram.com}.  The NFHS-5 is a clustered multistage sample survey that interviewed over 635,000 households and over 720,000 women between the ages of 15 and 49.

If we were to compute the prevalence of malnutrition in the prepregnancy period simply by estimating how many nonpregnant women are underweight, without statistical adjustment, we would produce a biased estimate.  Not all nonpregnant women will conceive in the future, and some are far more likely than others to conceive soon.  Further, women who are unlikely to conceive soon have systematically different nutritional status than those who are more likely to conceive soon. The bias generated by equating nonpregnant status with prepregnant status is especially important in India, where social hierarchies based on age and sex mean that women in their childbearing years are particularly vulnerable to undernutrition \citep{gupta1995life, coffey2021neonatal}.  

Therefore, we estimate the prevalence of prepregnancy underweight constructing a nonparametric reweighting function based on observable characteristics of pregnant women, and applying it to the body mass index scores of nonpregnant women.  This approach, which was developed by \cite{coffey2015prepregnancy}, approximates a synthetic cohort of prepregnant women.  Births outside of marriage are extremely rare -- more than 99\% of pregnant women in the NFHS-5 are married at the time of the survey -- so we restrict our sample to married women. Additionally, it is extremely rare (less than 1\%) for pregnant women to report using a method of modern contraception in the month before conception. Thus, we also restrict our sample to nonpregnant women currently not using modern contraception, and pregnant women who weren't using modern contraception in the month before conception. 

\subsection{Key Variables}

{\bf{Body mass index.}} Body mass index (BMI) is a person's weight in kilograms divided by her height in meters squared.  The WHO classifies someone as ``underweight'' if her body mass index is less than 18.5 kg/m$^2$.  The NFHS-5 measured women's weights and heights at the time of the survey.  BMI is missing for about 3\% of married women; these observations are dropped from the analysis.

{\bf{Pregnancy.}}  The NFHS-5 asks women ages 15 to 49 whether they are currently pregnant.  We drop women from the analysis who report pregnancies of one or two months duration because they are a select group.  Not all women are aware of early pregnancies and willing to report them.\footnote{Table \ref{regression3plusmopreg} in the online appendix shows that, among other differences, women who report being one or two months pregnant are more likely to be in the upper half of the wealth distribution than women who report three or more months of pregnancy. Table \ref{gestdur_distribution} in the online appendix shows that a smaller proportion of pregnant women report one or two months of pregnancy relative to months three through eight.}  We compute gestational duration using the reported number of months since the start of the woman's last menstrual period.  If that question was not answered, we use the woman's answer to the question, ``How many months pregnant are you?'' to measure gestational duration.\footnote{This substitution is made for 5\% of the pregnant sample.}  

{\bf{Social group.}} The NFHS asks respondents whether they belong a Scheduled Caste, a Scheduled Tribe, or an Other Backward Class.  Separately, it asks respondents their religion, listing Hindu, Muslim, Christian, Sikh, Buddhist, Jain, and ``others'' as options. 

Following the IHDS, which emphasizes the study of social inequality, we combine the caste and religion variables into five mutually exclusive groups \citep{IHDS2005}.  We code as Adivasi or Dalit those women who report belonging to a Scheduled Tribe or a Scheduled Caste, respectively, regardless of religion.  Women who report their religion as Muslim, and who do not belong to a Scheduled Caste or Tribe, are coded as Muslim. Hindu or Sikh women belonging to an Other Backward Class are coded as OBC.  Hindu women who do not belong to a Scheduled Tribe, a Scheduled Caste, or an Other Backward Class are coded as Forward caste Hindu.  2.3\% percent of married women interviewed by the NFHS-5 are dropped from the analysis because do not fall into any of these five categories.  They may identify as Christian, Jain, Sikh, or as having no religious affiliation.   

\subsection{Estimation of prepregnancy nutrition}

We apply a nonparametric reweighting function to the body mass index scores of nonpregnant women to make their distribution of the observable characteristics match that of pregnant women. This counterfactual distribution up-weights nonpregnant women who are more likely to become pregnant and down-weights those who are less likely to become pregnant. This method addresses the challenge that we cannot directly observe the nutrition status of women at the time of conception. 
 
The reweighting function \(\psi\) is defined as \(\psi(x)=[f(x|P=1)]/[f(x|P=0)]\), the ratio of the weighted counts of pregnant to nonpregnant women in each bin \(x\).  A bin is a unique combination of categories of the reweighting variables. To estimate the prevalence of prepregnancy underweight, the reweighting function is multiplied by each woman's sampling weight: \(\tilde{W}_i = W_i * \psi(x_i)\), where nonpregnant woman \(i\) in bin \(x_i\) has a sampling weight \(\tilde{W}_i\). A key assumption is that, conditional on the binning variables, the distribution of BMI among nonpregnant women matches that of women who are beginning pregnancy.

We construct 2,560 unique bins from the intersection of the categorical variables described below.  For pregnant women, variables are coded according to what was true when the woman became pregnant.  For nonpregnant women, variables are constructed according to what was true at the time of the survey.  These variables were chosen based on the literature about determinants of fertility in India \citep{jha2025effects, singh2022key, dharmalingam2014determinants}.  Table \ref{tab:regression} in the online appendix presents the results of OLS linear probability models showing that these variables predict pregnancy in the NFHS-5.  Models are estimated separately for each social group, and for all five groups combined. 

The variables used in the nonparametric reweighting technique are:
\begin{itemize}
% \item {\bf{not using modern contraception:}} We construct a binary variable that, for nonpregnant women, indicates that they are not using modern method of contraception in the month of the survey.  For pregnant women, it indicates that they did not use modern contraception in the month before the start of the pregnancy.  For pregnant women, the information is obtained from the reproductive calendar, which offers a month by month history of the last 6 years of a woman’s experiences with childbearing and contraceptive use. \footnote{We follow the NFHS-5's definition of modern contraception.  The NFHS-5 counts the following methods as modern contraception: pill, IUD, injections, diaphragm, condom (male or female), sterilization (male or female), lactational amenorrhea, foam or jelly, emergency contraception, and standard days method. Less than 1\% of pregnant women report having used a method of modern contraception in the month before they became pregnant.} 
\item {\bf{less than a primary education:}} This is a binary indicator for not having finished primary school.
\item {\bf{rural resident:}} A respondent is coded as a rural resident if she was interviewed in a rural primary sampling unit. 
\item {\bf{does not have a boy child:}} This is a binary indicator that a woman does not have any sons.
\item {\bf{age group:}} 15-19, 20-24, 25-29, 30-49 years old
\item {\bf{parity and birth spacing:}} We construct a categorical variable that captures both parity and birth spacing. Women who are pregnant with their first child are coded as parity 0, as are nonpregnant women who have no children.  For pregnant women, birth spacing is measured from the month of the previous live birth to the end of the current pregnancy, assuming that the current pregnancy will last nine months. For nonpregnant women, we assign the birth spacing that would be true if they became pregnant in the month of the survey.  Therefore, birth spacing is measured from the month of the previous live birth to nine months after the date of the interview. 
We construct the following categories: (1) Parity 0; (2) Parity 1, below 2 years spacing; (3) Parity 1, 2–3 years spacing; (4) Parity 1, 3+ years spacing; (5) Parity 2, 2 years spacing; (6) Parity 2, 2–3 years spacing; (7) Parity 2, 3+ years spacing; (8) Parity 3+, below 2 years spacing; (9) Parity 3+, 2–3 years spacing; and (10) Parity 3+, 3+ years spacing.
\item {\bf{household wealth quartile:}} The NFHS-5 publishes a wealth index score based on a principal components analysis of household asset ownership (for example, a bicycle, a television, or a car) and housing characteristics (for example, source of drinking water, toilet facilities, and flooring materials) \citep{iips_icf_nfhs5_2022}. We split the wealth index score into quartiles. 
\end{itemize}

Bins that contain nonpregnant women, but do not contain pregnant women, do not pose a problem for the analysis; they receive a weight of zero.  However, bins that contain pregnant women, but do not contain nonpregnant women, must be dropped from the analysis. Table \ref{tab:sampledropped} in the online appendix shows the sample sizes of pregnant women used in the subgroup analyses, and the percent of pregnant women dropped because they are in a bin that contains no nonpregnant women.  

The reweighting procedure was conducted within each social group.  The national estimate of prepregnancy underweight uses social group as a reweighting variable.  When we decompose differences in the prevalence of prepregnancy underweight by social group, we compute estimates of prepregnancy underweight within each social group's parity and birth spacing bins, and within each social group's wealth quartile bins.  

Standard errors are constructed using cluster bootstrapping.  We draw 30,125 clusters\footnote{This is slightly less than the number of primary sampling units in the NFHS-5 because we drop women from the analysis who do not belong to one of the social groups we study.} with replacement and compute the proportion of prepregnancy underweight in that sample, and repeat the procedure 1,000 times.  We present confidence intervals spanning the 5th and 95th percentiles of the distribution of resulting proportions of prepregnancy underweight.

\subsection{Kitagawa decomposition of differences in prepregnancy underweight between marginalized groups and Forward castes}

We conduct two Kitagawa decompositions \citep{kitagawa1955components}: the first decomposes gaps between marginalized groups and Forward caste Hindus using the parity and birth spacing categories, and the second decomposes gaps using wealth categories.  The overall gap in the prevalence of prepregnancy underweight between a marginalized group and Forward caste Hindus can be expressed as the sum of differences ``explained'' by the variable of interest and differences that remain ``unexplained:''
%HERE

\begin{equation}
\label{eq:kitagawa}
\Delta \equiv \bar Y^{g} - \bar Y^{f}
= \Delta_{explained} + \Delta_{unexplained}
\end{equation}

\begin{equation} \label{eq_explained}
    \Delta_{explained} = \sum_{j=1}^{J} \big(p_j^{g} - p_j^{f}\big)\,\frac{Y_j^{g} + Y_j^{f}}{2}
\end{equation}

\begin{equation} \label{eq_unexplained}
    \Delta_{unexplained} = \sum_{j=1}^{J} \big(Y_j^{g} - Y_j^{f}\big)\,\frac{p_j^{g} + p_j^{f}}{2}
\end{equation}

\begin{itemize}
    \item where \(j=1,...,J\) indexes categories of the decomposition variable (for parity and birth spacing J=10, for wealth quartiles J=4)
    \item \( \bar Y^g\) and \(\bar Y^f\) are the prevalence of prepregnancy underweight for groups \(g\) (marginalized group) and \(f\) (Forward caste Hindu)
    \item \(Y_j^g\) and \(Y_j^f\) are the means of prepregnancy underweight within category \(j\) of groups \(g\) and \(f\) respectively
    \item \(p_j^g\) and \(p_j^f\) are the proportions of women in each category \(j\) of the variable of interest
    
\end{itemize}

The percent of the gap explained by the variable of interest is computed as:

\begin{equation} \label{eq_pct}
    \frac{\Delta_{explained}}{\Delta_{total}}.
\end{equation}

\subsection{Statement on use of AI}\label{sec:AIuse}

We used OpenAI ChatGPT (GPT-5.2) to assist with Stata programming, table and figure formatting, and editing. Use cases included (i) requests to rewrite existing Stata code to improve runtime, (ii) requests to generate code in the Stata packages listtex and esttab to achieve desired table layouts, (iii) assistance with navigating Stata documentation to achieve desired figure layouts, and (iv) copyediting to reduce word count.  Sample prompts for these use cases can be found in Appendix section \ref{sec:aiprompts}.

\section{Results}

\subsection*{Descriptive statistics}

Table \ref{tab:sumstats} shows sample means for the variables we use in the reweighting decomposition, separately by social group and pregnancy status. Pregnant and nonpregnant women have different characteristics, on average. Pregnant women are younger and at lower parities: 43\% are aged 20–24 (versus 12\% of nonpregnant women), and 42\% have no prior births (versus 8\% of nonpregnant women). This highlights the potential biases of estimating prepregnancy underweight from the BMIs of nonpregnant women without adjustment for selection into pregnancy.

There are meaningful differences in these predictors of pregnancy across social groups. Nearly half of nonpregnant Adivasi women have less than primary education, compared to one in five nonpregnant Forward caste women. Over half of nonpregnant Adivasi women and over one fourth of Dalit women live in households in the bottom quartile of the wealth distribution, compared with only 10\% of nonpregnant Forward caste women. Adivasi and Dalit women are also overrepresented in the bottom quartile of household wealth among pregnant women.

\subsection*{Prepregnancy underweight by social group}

Figure \ref{fig:mainresult} shows estimates of the prevalence of prepregnancy underweight by social group, and for all five groups combined. With the exception of Forward caste Hindus and Muslims, confidence intervals do not overlap for any social groups, indicating statistically significant differences in the prevalence of prepregnancy underweight at the 95\% level. The burden of prepregnancy underweight is highest among Adivasi women (27\%), followed by Dalit women (23\%), and OBC (20\%) women.  Forward caste Hindus (16\%) and Muslims (14\%) have a lower prevalence of prepregnancy underweight. Combining all five social groups, we find that 20\% of prepregnant women are underweight.\footnote{Our sample drops a small number of women who do not fall into one of the five social groups used in our analysis.  If we estimate prepregnancy underweight using the full NFHS-5 sample, we find that the same proportion of prepregnant women are underweight (19.83\%).} Table \ref{specifications} in the online appendix shows that the results are robust to changes in the covariates used in the reweighting procedure. 

Although not the focus of this paper, we observe that Adivasi, Dalit, and OBC women do not compensate for prepregnancy underweight with higher weight gain during pregnancy.\footnote{The US Institute of Medicine (IOM) recommends that women who begin pregnancy underweight gain 28 to 40 pounds during pregnancy \citep{iom2009weightgain}. Women who have normal or overweight BMIs at the beginning of pregnancy should gain less.  We cite the IOM guidelines, which have been endorsed by the United Kingdom's National Institute for Health and Care Excellence \citep{NICE2025}, because the Government of India has not issued guidelines on weight gain in pregnancy.}  Table \ref{tab:addnut} in the online appendix uses NFHS-5 data to show estimates of weight gain during pregnancy for each social group. We find that no group's gestational weight gain exceeds 20 pounds, and that there are no statistically significant differences in gestational weight gain across social groups. The methods for deriving these estimates were developed by \cite{coffey2015prepregnancy} and are detailed in Section \ref{sec:weightgain} of the online appendix.\footnote{Table \ref{tab:addnut} also computes the prevalence of prepregnancy obesity, defined as having a BMI over 30, by social group and for India as a whole, using the same nonparametric method we used to compute the prevalence of prepregnancy underweight.  There are small differences by social group: 11\% of Forward caste Hindu and Muslim women are obese before pregnancy, compared to 5\% of Adivasi women, 7\% of Dalit women, and 8\% of OBC women.  We include these estimates for completeness and to encourage future research on maternal nutrition.  Although obesity rates are increasing in India \citep{luhar2020forecasting}, obesity among women who are likely to become pregnant is relatively uncommon: 8\% of Indian women are estimated to be obese before pregnancy.} 


\subsection*{Social group differences in parity, birth spacing, and household wealth}

In order to investigate explanations for these gaps, panels A, C, and E of Figure \ref{fig:sixpanel} show the distributions of pregnant women, by social group, into parity, birth spacing, and household wealth categories.  Panels B, D, and F of \ref{fig:sixpanel} show estimates of the prevalence of prepregnancy underweight, for all social groups combined, within parity, birth spacing, and household wealth categories.  

\subsubsection*{Parity}

Panel A of Figure \ref{fig:sixpanel} shows that Forward caste Hindu pregnant women are more likely than pregnant women from other social groups to be pregnant with their first child (parity 0). Almost half of pregnant Forward caste Hindu women are at parity 0, compared to 38-39\% of Muslim and Adivasi pregnant women. Muslim pregnant women are more likely than women from other social groups to be at higher parities, with 16\% at parity 3 or higher compared to 5\% among Forward caste Hindu women and 8\% among OBC women. Although present, fertility differences across social groups are not very large.\footnote{Table \ref{tab:tfr} in the online appendix complements the information in Panel A of Figure \ref{fig:sixpanel} by showing total fertility rates (TFR) for each of the five social groups based on births in the three years prior to the NFHS-5.  TFR for Forward caste Hindu women is 1.6, compared to 2.3 among Muslims, and 2.1 among Adivasis and Dalits.}  Fertility in India today is very close to replacement; higher order births are becoming rare.

Panel B of Figure \ref{fig:sixpanel} shows no statistically significant differences in prepregnancy underweight by parity. This finding does not rule out that as women progress through their childbearing careers, they may be less likely to be underweight.  The cross-sectional nature of the data limits what we can conclude about the relationship between parity progression and maternal nutrition at the individual level.  Poorer women are more likely to be undernourished at any parity, and are more likely to progress to higher parities than wealthier women.  

\subsubsection*{Birth spacing}

Panel C of Figure \ref{fig:sixpanel} summarizes birth spacing for the 14,240 pregnant women in the NFHS-5 who had at least one prior birth.  Compared to women from other social groups, Dalit women were most likely to have birth spacing below two years.  28\% of Dalit pregnant women, compared to about 21\% of Forward caste Hindu, Adivasi, and Muslim pregnant women, are expected to give birth within two years of a prior birth.  Three or more years of birth spacing is modal for all social groups.  55\% of Muslim pregnant women are expected to have at least three years birth spacing, compared to 40\% of Dalit pregnant women. Panel D of Figure \ref{fig:sixpanel} shows that the estimated prevalence of prepregnancy underweight by birth spacing category: women with 3 or more years spacing from the prior birth are less likely to be underweight than those with less than 3 years birth spacing.

\subsubsection*{Household Wealth}

Panel E of Figure \ref{fig:sixpanel} shows stark differences in the distribution of household wealth across Indian social groups. Half of Adivasi pregnant women live in households in the lowest wealth quartile, compared to only 12\% of Forward caste Hindu pregnant women. By contrast, nearly 43\% of Forward caste Hindu pregnant women are in the top quartile, compared with fewer than one in five Dalit or Adivasi pregnant women. 

Panel F of Figure \ref{fig:sixpanel} estimates prepregnancy underweight by household wealth category.  In contrast to panels B and D, which show that the prevalence of prepregnancy underweight varies little by parity or by birth spacing, a clear and quantitatively important pattern emerges in panel F.  The prevalence of prepregnancy underweight is inversely related to household wealth.  Almost thirty percent of women in the poorest quartile are estimated to be underweight at the start of pregnancy, compared to 11\% in the top quartile.  

\subsection*{Decompositions of differences in prepregnancy underweight across social groups}\label{sec:kitagawa}

Table \ref{decomposition} summarizes Kitagawa decompositions of gaps in the prevalence of prepregnancy underweight between women from marginalized groups and Forward caste Hindu women.  The top row of Table \ref{decomposition} shows the raw gaps from Figure \ref{fig:mainresult}, the largest of which is the Adivasi-Forward gap of 10.8 percentage points.  Panel A shows that the parity and birth spacing categories explain little of the Forward caste advantage in prepregnancy underweight relative to Adivasi, Dalit, and OBC women (3\% to 11\% of the gaps). 

Differences in the distributions of household wealth, in contrast, explain a large portion of the Forward caste Hindu advantage in prepregnancy underweight relative to Adivasis, Dalits, and OBCs: 68\%, 62\%, and 59\% of the respective gaps.  The decomposition of the gap between Muslims and Forward castes suggests that if Muslims had the same distribution of wealth as Forward caste Hindus, they might face a lower burden of prepregnancy underweight than Forward caste Hindus. 

\section{Discussion}

This paper contributes the most recent estimate of prepregnancy underweight for India and the first estimates stratified by social group. We find that in 2019-2021, about one in five women in India was underweight at the beginning of pregnancy. For comparison, fewer than 5\% of women in the United States are underweight at the start of pregnancy \citep{zong2022maternal}.  Nevertheless, this represents significant progress in maternal health because the prevalence of prepregnancy underweight has come down by over 50\% since 2005.  

Within India, the burden of maternal undernutrition falls disproportionately on Adivasi, Dalit, and OBC women, who have a prevalence of prepregnancy underweight 11, 7, and 5 percentage points higher than forward caste Hindus, respectively.  The prevalence of prepregnancy underweight among Forward caste Hindu and Muslim women is similar; this finding coheres with other research that finds advantages in Muslim women's and children's health relative to Hindus'.   

The Kitagawa decompositions show that higher fertility and more closely spaced births among Adivasis, Dalits, and OBCs do not explain why they are more likely to be underweight before pregnancy than Forward caste Hindu women.  This is in part because the differences by social group in parity and birth spacing are small, and in part because these variables are only weakly correlated with prepregnancy underweight in the most recent Indian data. Differences in household wealth, in contrast, do explain large proportions of the gaps.  Future studies might investigate whether this is because wealthier households are able to afford food, because wealthier households can afford labor saving devices such as vehicles or washing machines, or because wealthier households offer better protection against disease. 

A limitation of this study is that a meaningful fraction of the gaps in the prevalence of prepregnancy underweight remains unexplained after accounting for differences in household wealth.  Could remaining gaps in the prevalance of prepregnancy underweight be explained by differences in intrashousehold allocation of food across caste groups? The NFHS does not have information about the relative quantities of food consumed by different household members that might help us answer this question.  However, the literature and some evidence from the IHDS suggest that caste differences in intrahousehold food allocation may not explain gaps in the prevalence of prepregnany underweight.  \cite{harris2017determinants}'s review  of intrahousehold food allocation among adults in India suggests that the opposite may be true: oppressed caste households may distribute food more equally between men and women than Forward caste Hindu households because oppressed caste women are more likely to work outside the home than dominant caste women \citep{chakrabarty1996gender, eswaran2013status}.  To provide some preliminary evidence on the relationship between caste and intrahousehold food allocation, we analyzed an IHDS (2011) question which asked ever-married women whether they eat after or with the men in their households.  Figure \ref{fig:ihds_eat_last} in the online appendix shows the proportion of women who report eating after men is very similar across social groups.\footnote{\cite{desai2014muslim} analyzed IHDS (2005) data and found that Muslim women \emph{were} less likely to report eating after the men in their household than Hindu women.}  

The literature on racial differences in birth outcomes in the United States increasingly points to racism-related stress among mothers as an important reason for higher rates of preterm birth and lower birth weights among Black infants \citep{hogue2005stress, sosnaud2021cross}. Could differences in discrimination-related stress explain the remaining caste gaps in maternal undernutrition? A limitation of this study is that, to our knowledge, no population-level data on biomarkers such as cortisol and C-reactive protein exist that could help answer this question. 

Despite these limitations, our research nevertheless has implications for policy and research.  The finding that maternal undernutrition remains high at the population level underscores the importance of the 2013 National Food Security Act's maternity entitlements to pregnant women. The law promised nearly all pregnant women a Rs. 6,000 benefit, but in practice, only a fraction of the funds that would be needed for a universal program have been allocated.\footnote{In 2020, allocation to maternity entitlements was only enough to cover about 1 in 7 women who gave birth \cite{dreze2021derailed}.  Maternity entitlements are often restricted to the first child and are conditional on the woman receiving care from government health clinics \citep{kalra2020birth}.}  If fully funded, indexed for inflation, and given at the start of pregnancy, cash transfers during pregnancy could help address India's remaining burden of prepregnancy underweight \citep{coffey2016underweight}. 

Programs and policies to address maternal undernutrition should focus not only on pregnanct women, but also on adolescents and not-yet-married women.  The time  between marriage and conception is short: on average, women who had their first births in the three years before the NFHS-5 were married for only 17.23 months before conception.  Cash or in-kind transfers to adolescent girls and young women in their natal households may improve prepregnancy health.  

Finally, this study highlights the value of representative, longitudinal monitoring data on health in pregnancy.  The cross sectional data that we analyze permit the construction of synthetic cohorts that yield valuable insights into maternal nutrition, and social group differences in maternal nutrition.  Yet, they do not allow for individual level analyses.  For example, we do not know what fraction of women who begin pregnancy underweight compensate for low prepregnancy body mass with adequate weight gain.  Nor do we know how many women who begin pregnancy underweight go on to have a preterm or low birthweight infant.  Investments in longitudinal health data will improve health policy for both mothers and children.


\newpage

\section*{Figures and Tables}

% \subsection*{Prepregnancy underweight by social group (Figure 1)}
\begin{figure}[H]
    \centering
    \includegraphics[width=\textwidth]
    {figures/figure1 prepregnancy underweight by subgroup.png}
    \caption{: Prepregnancy underweight by social group}
    \label{fig:mainresult}
    \parbox{1\linewidth}{\footnotesize NFHS-5 data. Estimates are generated using a nonparametric reweighting function applied to married nonpregnant women who don't report currently using modern contraception in the NFHS-5. Confidence intervals are computed using cluster bootstrapping.}
\end{figure}


\clearpage
\thispagestyle{empty}

\begin{landscape}
% \subsection*{Prepregnancy underweight by parity, birth spacing, and wealth quartile (Figure 2)}

\begin{figure}[p]
  \vspace*{-1cm}
  \makebox[\textwidth][c]{%
    % \includegraphics[width=0.95\textwidth]{figures/six panel figure.png}

    \includegraphics[width=1.2\textwidth]{figures/figure2 predictors by social group and gradients.pdf}
  }
  \caption{: Differences in parity, birth spacing, and wealth quartile by social group, and the relationship between these characteristics and prepregnancy underweight}
  \label{fig:sixpanel}

  \parbox{1\linewidth}{\footnotesize NFHS-5 data.}

\end{figure}

\end{landscape}



\begin{landscape}

\begin{table}[!htbp]
    \centering
    \begin{threeparttable}[t]
        \caption{:  Descriptive statistics by social group and pregnancy status}
        \label{tab:sumstats}

        \scriptsize
        \setlength{\tabcolsep}{2pt}
        \renewcommand{\arraystretch}{1.2}

        \begin{tabular}{l*{6}{>{\centering\arraybackslash}p{1.2cm}}@{\hspace{3em}}*{6}{>{\centering\arraybackslash}p{1.2cm}}}
\toprule
& \multicolumn{6}{c}{Pregnant women (3+ months)} & \multicolumn{6}{c}{Nonpregnant women} \\\\
Social Group & \tiny Adivasi & \tiny Dalit & \tiny OBC & \tiny Forward & \tiny Muslim & \tiny \shortstack{All five \\\\ social groups} & \tiny Adivasi & \tiny Dalit & \tiny OBC & \tiny Forward & \tiny Muslim & \tiny \shortstack{All five \\\\ social groups} \\\\
\midrule
\textbf{Binary Predictors of Pregnancy and Underweight}&.9674813&.9769595&.9805262&.9738441&.9678563&.9754397&.9839144&.9843819&.976019&1&.9876108&.985052\\
\hspace*{2em}not using modern contraception&.97&.98&.98&.97&.97&.98&.98&.98&.98&.00&.99&.99\\
\hspace*{2em}less than primary education&.36&.32&.24&.12&.35&.27&.32&.32&.35&.23&.27&.30\\
\hspace*{2em}rural resident&.89&.78&.76&.64&.66&.74&.91&.71&.80&.78&.60&.75\\
\hspace*{2em}does not have boy child&.66&.66&.71&.77&.59&.68&.64&.77&.58&.67&.60&.64\\
&.657608&.6629888&.7120327&.7709119&.5934445&.6848371&.6372807&.7669369&.5752585&.6720339&.5979783&.6399697\\
\textbf{Age Categories}&.657608&.6629888&.7120327&.7709119&.5934445&.6848371&.6372807&.7669369&.5752585&.6720339&.5979783&.6399697\\
\hspace*{2em}15 to 19&.13&.13&.11&.10&.12&.12&.07&.11&.10&.06&.13&.10\\
\hspace*{2em}20 to 24&.39&.45&.45&.39&.40&.42&.34&.36&.40&.31&.34&.36\\
\hspace*{2em}25 to 29&.30&.29&.31&.36&.29&.31&.37&.36&.30&.43&.26&.33\\
\hspace*{2em}30 to 49&.17&.13&.12&.16&.19&.15&.21&.17&.20&.20&.27&.21\\
&.1712409&.1272503&.1231072&.159552&.188175&.1465245&.2132325&.1728706&.1996052&.2033579&.2673149&.2056437\\
\textbf{Categories for parity and spacing from last birth}&.1712409&.1272503&.1231072&.159552&.188175&.1465245&.2132325&.1728706&.1996052&.2033579&.2673149&.2056437\\
\hspace*{2em}No births&.39&.41&.43&.49&.37&.42&.32&.40&.38&.36&.39&.37\\
\hspace*{2em}1 birth, \textless{}2y spacing&.06&.09&.08&.07&.05&.08&.02&.05&.08&.02&.02&.04\\
\hspace*{2em}1 birth, 2–3y spacing&.10&.10&.10&.08&.06&.09&.14&.05&.08&.09&.09&.08\\
\hspace*{2em}1 birth, \textgreater{}3y spacing&.17&.14&.16&.21&.15&.16&.22&.25&.15&.43&.23&.25\\
\hspace*{2em}2 births, \textless{}2y spacing&.03&.04&.03&.02&.03&.03&.00&.00&.05&.00&.02&.02\\
\hspace*{2em}2 births, 2–3y spacing&.05&.05&.05&.03&.04&.04&.04&.05&.00&.03&.02&.03\\
\hspace*{2em}2 births, \textgreater{}3y spacing&.03&.03&.02&.01&.04&.03&.00&.01&.01&.00&.00&.01\\
\hspace*{2em}3+ births, \textless{}2y spacing&.04&.04&.03&.02&.05&.04&.09&.04&.04&.01&.00&.04\\
\hspace*{2em}3+ births, 2–3y spacing&.08&.06&.06&.05&.12&.07&.11&.07&.10&.02&.13&.09\\
\hspace*{2em}3+ births, \textgreater{}3y spacing&.05&.04&.03&.02&.07&.04&.07&.07&.10&.03&.11&.08\\
&.0520137&.0388658&.0336344&.0193346&.072835&.0409029&.0736299&.0727488&.1009637&.0252251&.1060868&.0788187\\
\textbf{Wealth Categories}&.0520137&.0388658&.0336344&.0193346&.072835&.0409029&.0736299&.0727488&.1009637&.0252251&.1060868&.0788187\\
\hspace*{2em}1st (bottom) quartile&.51&.35&.23&.12&.27&.27&.42&.34&.26&.08&.26&.27\\
\hspace*{2em}2nd quartile&.25&.26&.25&.18&.22&.24&.27&.20&.29&.31&.17&.25\\
\hspace*{2em}3rd quartile&.16&.23&.26&.28&.25&.24&.21&.22&.23&.35&.11&.22\\
\hspace*{2em}4th quartile&.09&.16&.26&.43&.25&.25&.10&.24&.22&.25&.46&.27\\
&.0879128&.1576316&.260832&.4281044&.2534508&.246176&.0954724&.2419103&.2226783&.2491975&.4575896&.2685206\\
&.0879128&.1576316&.260832&.4281044&.2534508&.246176&.0954724&.2419103&.2226783&.2491975&.4575896&.2685206\\
\textbf{N}&2.26e+09&5.12e+09&7.96e+09&3.05e+09&3.67e+09&2.23e+10&8.93e+07&1.40e+08&1.76e+08&8.97e+07&1.12e+08&6.18e+08\\
\bottomrule
\end{tabular}


        \begin{tablenotes}
            \item NFHS-5 data. Sample means for women who are 3+ months pregnant are calculated based on what was true in the month of conception. Sample means for nonpregnant women (unmarried, not currently using modern contraception) are calculated based on what was true at the time of the survey.
        \end{tablenotes}
    \end{threeparttable}
\end{table}
\end{landscape}

\begin{landscape}

\begin{table}[H]
    \centering

    \begin{threeparttable}
        
    
    \small % shrink text
    \caption{: Decomposition of social group differences in the prevalence of prepregnancy underweight}
    \label{decomposition}
    % \setlength{\tabcolsep}{1pt}
    \renewcommand{\arraystretch}{1.5}
    % {\begin{tabular}{l*{4}{c}}
\toprule
            &\multicolumn{1}{c}{Adivasi-Forward}&\multicolumn{1}{c}{Dalit-Forward}&\multicolumn{1}{c}{OBC-Forward}&\multicolumn{1}{c}{Muslim-Forward}\\
\midrule
\midrule
percentage point difference in pre-pregnancy underweight&        5.07&        7.03&       10.78&       -1.65\\
            &            &            &            &            \\
\textbf{Panel A: Decompositions of parity + birthspacing}&            &            &            &            \\
pp difference within parity + birthspacing category&        4.54&        6.30&       10.51&       -1.48\\
pp different across parity + birthspacing category&        0.50&        0.72&        0.25&       -0.17\\
\% explained by parity + birthspacing&        9.84&       10.19&        2.30&       10.50\\
            &            &            &            &            \\
\textbf{Panel B. Decompositions of wealth}&            &            &            &            \\
pp difference within wealth category&        2.47&        2.94&        3.96&       -3.94\\
pp different across wealth category&        2.56&        4.07&        6.77&        2.25\\
\% explained by wealth&       50.61&       57.91&       62.75&     -136.45\\
\bottomrule
\end{tabular}
}
    {\begin{tabular}{l*{4}{c}}
\toprule
            &\multicolumn{1}{c}{Adivasi-Forward}&\multicolumn{1}{c}{Dalit-Forward}&\multicolumn{1}{c}{OBC-Forward}&\multicolumn{1}{c}{Muslim-Forward}\\
\midrule
\midrule
percentage point difference in prepregnancy underweight&        10.1&         7.2&         5.0&        -1.4\\
            &            &            &            &            \\
\textbf{Panel A: Decompositions of parity + birthspacing}&            &            &            &            \\
pp difference within parity + birthspacing category$^{1}$&        10.3&         6.8&         4.7&        -0.7\\
pp difference across parity + birthspacing category$^{2}$&        -0.1&         0.4&         0.2&        -0.8\\
\% explained by parity + birthspacing$^{3}$&        -1.3&         5.2&         4.4&        52.6\\
            &            &            &            &            \\
\textbf{Panel B. Decompositions of wealth}&            &            &            &            \\
pp difference within wealth category$^{1}$&         2.6&         2.6&         2.3&        -4.0\\
pp difference across wealth category$^{2}$&         7.5&         4.6&         2.7&         2.5\\
\% explained by wealth$^{3}$&        73.7&        63.7&        53.6&      -175.1\\
\bottomrule
\end{tabular}
}

    \begin{tablenotes}
        \item NFHS-5 data.  The ``percentage point difference in prepregnancy underweight'' is computed by subtracting the prevalence of prepregnancy underweight (from Figure \ref{fig:mainresult}) of Forward caste Hindus from that of each marginalized group. The ``\%s explained'' are computed using \cite{kitagawa1955components} decomposition. 
    \end{tablenotes}

    \end{threeparttable}
\end{table}

\end{landscape}




% % % % % % % % % % % % % 


\bibliographystyle{apalike}
\bibliography{writing/references_apalike_apa_braced}

\newpage
\clearpage
\appendix

\section*{Appendix}



% Reset and format appendix numbering
\setcounter{section}{-1}
\renewcommand{\thesection}{S.\arabic{section}}
\renewcommand{\thesubsection}{S.\arabic{section}.\arabic{subsection}}

% Make tables/figures labeled A.1, A.2, etc.
\setcounter{table}{0}
\renewcommand{\thetable}{A\arabic{table}}
\setcounter{figure}{0}
\renewcommand{\thefigure}{A\arabic{figure}}

% Now numbered sections start from S.1
\stepcounter{section}

\section{S.\arabic{section} Standalone figures and tables} \label{sec:appendixexhibits}

% \newgeometry{left=1.2in, right=1.2in}  % or smaller margins

\begin{figure}[H]
    \centering
    \includegraphics[width=\textwidth]
    {figures/figureA1 women eating last ihds.png}
    \caption{: Proportion of women who report eating last, by social group}
    \label{fig:ihds_eat_last}
    \parbox{1\linewidth}{\footnotesize IHDS (2011) data. Sample includes 34,216 ever-married women ages 15 to 49.  We apply the same social group definitions to the IHDS data as we do to the NFHS data.}
\end{figure}

\clearpage

\begin{table}[H]

\begin{threeparttable}
    
    \centering
    \scriptsize % shrink text
    
    \caption{: Women reporting one or two months of pregnancy are different than those reporting three or more months of pregnancy}
    \label{regression3plusmopreg}
    \setlength{\tabcolsep}{1pt}
    \renewcommand{\arraystretch}{0.81}
    {{
\def\sym#1{\ifmmode^{#1}\else\(^{#1}\)\fi}
\begin{tabular}{l*{1}{c}}
\toprule
                    &\multicolumn{1}{c}{\shortstack{reports of 1 or 2 months of pregnancy}}\\
\midrule
\textbf{Binary predictors}&                     \\
\addlinespace
\hspace*{1em}Not using modern contraception&       0.045\sym{***}\\
                    &    (0.0110)         \\
\addlinespace
\hspace*{1em}Less than primary education&       0.011\sym{*}  \\
                    &    (0.0050)         \\
\addlinespace
\hspace*{1em}Rural resident      &       0.009         \\
                    &    (0.0056)         \\
\addlinespace
\hspace*{1em}Does not have boy child&       0.006         \\
                    &    (0.0051)         \\
\addlinespace
\textbf{Age categories} \\ (15–19 omitted)&                     \\
\addlinespace
\hspace*{1em}20–24               &      -0.012         \\
                    &    (0.0073)         \\
\addlinespace
\hspace*{1em}25–29               &      -0.061\sym{***}\\
                    &    (0.0080)         \\
\addlinespace
\hspace*{1em}30–49               &      -0.099\sym{***}\\
                    &    (0.0094)         \\
\addlinespace
\textbf{Parity \& time since last live birth categories} \\ (No prior births omitted)&                     \\
\addlinespace
\hspace*{1em}1 birth, below 2y spacing&      -0.129\sym{***}\\
                    &    (0.0040)         \\
\addlinespace
\hspace*{1em}1 birth, 2–3y spacing&      -0.121\sym{***}\\
                    &    (0.0043)         \\
\addlinespace
\hspace*{1em}1 birth, above 3y spacing&       0.134\sym{***}\\
                    &    (0.0076)         \\
\addlinespace
\hspace*{1em}2 births, below 2y spacing&      -0.116\sym{***}\\
                    &    (0.0049)         \\
\addlinespace
\hspace*{1em}2 births, 2–3y spacing&      -0.107\sym{***}\\
                    &    (0.0052)         \\
\addlinespace
\hspace*{1em}2 births, above 3y spacing&       0.137\sym{***}\\
                    &    (0.0100)         \\
\addlinespace
\hspace*{1em}3+ births, below 2y spacing&      -0.088\sym{***}\\
                    &    (0.0065)         \\
\addlinespace
\hspace*{1em}3+ births, 2–3y spacing&      -0.081\sym{***}\\
                    &    (0.0066)         \\
\addlinespace
\hspace*{1em}3+ births, above 3y spacing&       0.179\sym{***}\\
                    &    (0.0132)         \\
\addlinespace
\textbf{Wealth quartiles} \\ (1st quartile omitted)&                     \\
\addlinespace
\hspace*{1em}2nd quartile        &       0.008         \\
                    &    (0.0054)         \\
\addlinespace
\hspace*{1em}3rd quartile        &       0.010         \\
                    &    (0.0058)         \\
\addlinespace
\hspace*{1em}4th quartile        &       0.026\sym{***}\\
                    &    (0.0069)         \\
\addlinespace
\hspace*{1em}Constant            &       0.087\sym{***}\\
                    &    (0.0150)         \\
\midrule
\textbf{N}          &      28,114         \\
\bottomrule
\end{tabular}
}
}
    \begin{tablenotes}
        \item Observations are women interviewed by the NFHS-5 who are in one of the five social groups we study and who report being currently pregnant. Standard errors, shown in parentheses, are clustered at the primary sampling unit level. 
        \item * $p<0.5$, ** $p<0.01$
    \end{tablenotes}
\end{threeparttable}
\end{table}

% \restoregeometry



\newpage

\begin{table}[H]
    \centering

    \begin{threeparttable}
            
    \caption{: Distribution of gestational ages suggests under-reporting of early pregnancy}
    \label{gestdur_distribution}
    \small
    \setlength{\tabcolsep}{5pt}
    
        \begin{tabular}{lcc} \toprule Month of gestation & Number of women & Percent of pregnant women \\ \midrule
\midrule
1                   &       1,308&        4.63\\
2                   &       2,442&        8.64\\
3                   &       3,502&       12.38\\
4                   &       3,602&       12.74\\
5                   &       3,458&       12.23\\
6                   &       3,506&       12.40\\
7                   &       3,400&       12.02\\
8                   &       3,692&       13.06\\
9                   &       2,679&        9.47\\
10                  &         611&        2.16\\
11                  &          78&        0.28\\
Total               &      28,280&      100.00\\
\bottomrule \end{tabular}

    
    \small
    \begin{tablenotes}
        \item  
    \end{tablenotes}

    \begin{tablenotes}
    \footnotesize
    \item 
        Observations are women interviewed by the NFHS-5 who are in one of the five social groups we study who report being currently pregnant.
    \end{tablenotes}
    
    \end{threeparttable}
\end{table}

%\setcounter{table}{1}
%\renewcommand{\thetable}{A\arabic{table}}
\newpage
% \begin{landscape}
% \subsection{Sample sizes and percent dropped by subgroup}

\newpage
\begin{table}[!htbp]
    \centering
    \begin{threeparttable}
        \caption{: Variables used in the nonparametric reweighting predict pregnancy}
        \label{tab:regression}

        \scriptsize
        \setlength{\tabcolsep}{1pt}
        \renewcommand{\arraystretch}{0.81}

        % {{
\def\sym#1{\ifmmode^{#1}\else\(^{#1}\)\fi}
\begin{tabular}{l*{6}{c}}
\toprule
                    &\multicolumn{1}{c}{Adivasi}&\multicolumn{1}{c}{Dalit}&\multicolumn{1}{c}{OBC}&\multicolumn{1}{c}{Muslim}&\multicolumn{1}{c}{Forward}&\multicolumn{1}{c}{All 5 Social Groups}\\
\midrule
\textbf{Binary predictors of pregnancy and underweight}&                     &                     &                     &                     &                     &                     \\
\addlinespace
Uses modern contraception&      -.0752\sym{***}&      -.0727\sym{***}&      -.0641\sym{***}&      -.0782\sym{***}&      -.0512\sym{***}&      -.0671\sym{***}\\
                    &    (.00308)         &    (.00191)         &    (.00144)         &    (.00259)         &    (.00184)         &   (.000907)         \\
\addlinespace
None or incomplete primary education&     -.00853\sym{***}&      .00111         &      .00363\sym{***}&      .00155         &    -.000858         &      .00157\sym{*}  \\
                    &    (.00269)         &    (.00178)         &    (.00124)         &    (.00273)         &    (.00166)         &   (.000808)         \\
\addlinespace
Rural resident      &       .0123\sym{*}  &      .00582\sym{***}&      .00299\sym{*}  &       .0056\sym{*}  &      .00721\sym{***}&      .00486\sym{***}\\
                    &    (.00707)         &    (.00221)         &    (.00157)         &    (.00305)         &    (.00202)         &    (.00101)         \\
\addlinespace
Has boy child       &      -.0219\sym{***}&       -.016\sym{***}&      -.0252\sym{***}&      -.0132\sym{***}&      -.0208\sym{***}&      -.0206\sym{***}\\
                    &    (.00475)         &    (.00273)         &    (.00206)         &    (.00427)         &    (.00245)         &    (.00128)         \\
\addlinespace
\textbf{Age categories}&                     &                     &                     &                     &                     &                     \\
\addlinespace
20–24               &      -.0105         &      .00457         &       .0122         &       .0168         &       .0078         &      .00856\sym{*}  \\
                    &     (.0118)         &    (.00964)         &     (.0079)         &      (.012)         &     (.0133)         &    (.00467)         \\
\addlinespace
25–29               &      -.0371\sym{***}&      -.0327\sym{***}&      -.0294\sym{***}&      -.0244\sym{**} &      -.0242\sym{*}  &      -.0297\sym{***}\\
                    &     (.0119)         &     (.0097)         &    (.00787)         &     (.0119)         &     (.0132)         &    (.00465)         \\
\addlinespace
30–34               &      -.0769\sym{***}&      -.0696\sym{***}&      -.0706\sym{***}&      -.0748\sym{***}&      -.0626\sym{***}&      -.0709\sym{***}\\
                    &     (.0118)         &    (.00964)         &    (.00779)         &     (.0118)         &     (.0132)         &    (.00463)         \\
\addlinespace
\textbf{Parity \& time since last live birth categories}&                     &                     &                     &                     &                     &                     \\
\addlinespace
1 birth, <2y spacing&      -.0919\sym{***}&      -.0769\sym{***}&      -.0986\sym{***}&       -.145\sym{***}&        -.11\sym{***}&       -.101\sym{***}\\
                    &    (.00935)         &     (.0074)         &    (.00575)         &     (.0103)         &    (.00868)         &    (.00351)         \\
\addlinespace
1 birth, 2–3y spacing&      .00887         &      -.0339\sym{***}&      -.0407\sym{***}&      -.0866\sym{***}&      -.0793\sym{***}&      -.0483\sym{***}\\
                    &     (.0161)         &     (.0113)         &    (.00869)         &     (.0151)         &      (.011)         &    (.00516)         \\
\addlinespace
1 birth, 3+y spacing&      -.0524\sym{***}&      -.0843\sym{***}&      -.0847\sym{***}&      -.0898\sym{***}&      -.0982\sym{***}&      -.0852\sym{***}\\
                    &     (.0114)         &    (.00716)         &    (.00536)         &     (.0107)         &    (.00787)         &    (.00337)         \\
\addlinespace
2 births, <2y spacing&       -.101\sym{***}&       -.109\sym{***}&       -.123\sym{***}&       -.155\sym{***}&        -.13\sym{***}&       -.122\sym{***}\\
                    &    (.00904)         &    (.00707)         &    (.00536)         &     (.0101)         &    (.00811)         &    (.00333)         \\
\addlinespace
2 births, 2–3y spacing&      -.0603\sym{***}&      -.0835\sym{***}&      -.0975\sym{***}&       -.095\sym{***}&       -.114\sym{***}&      -.0924\sym{***}\\
                    &      (.011)         &    (.00874)         &    (.00615)         &     (.0124)         &    (.00898)         &    (.00392)         \\
\addlinespace
2 births, 3+y spacing&      -.0704\sym{***}&      -.0953\sym{***}&      -.0991\sym{***}&       -.109\sym{***}&       -.107\sym{***}&      -.0977\sym{***}\\
                    &     (.0085)         &    (.00649)         &    (.00494)         &    (.00975)         &    (.00753)         &    (.00308)         \\
\addlinespace
3+ births, <2y spacing&       -.077\sym{***}&      -.0981\sym{***}&       -.107\sym{***}&       -.114\sym{***}&       -.111\sym{***}&       -.101\sym{***}\\
                    &    (.00914)         &    (.00723)         &    (.00565)         &      (.011)         &    (.00898)         &    (.00358)         \\
\addlinespace
3+ births, 2–3y spacing&      -.0529\sym{***}&      -.0731\sym{***}&       -.086\sym{***}&         -.1\sym{***}&      -.0994\sym{***}&      -.0815\sym{***}\\
                    &    (.00994)         &    (.00817)         &    (.00641)         &     (.0113)         &     (.0101)         &    (.00389)         \\
\addlinespace
3+ births, 3+y spacing&        -.07\sym{***}&      -.0947\sym{***}&      -.0988\sym{***}&       -.121\sym{***}&       -.104\sym{***}&      -.0979\sym{***}\\
                    &    (.00828)         &    (.00655)         &    (.00491)         &    (.00943)         &    (.00756)         &    (.00308)         \\
\addlinespace
\textbf{Wealth categories}&                     &                     &                     &                     &                     &                     \\
\addlinespace
2nd quartile        &     -.00235         &      -.0077\sym{***}&     -.00659\sym{***}&     -.00774\sym{**} &     -.00191         &     -.00566\sym{***}\\
                    &    (.00262)         &    (.00222)         &     (.0018)         &    (.00391)         &    (.00309)         &    (.00112)         \\
\addlinespace
3rd quartile        &      .00293         &     -.00584\sym{**} &     -.00863\sym{***}&     -.00146         &      .00295         &     -.00433\sym{***}\\
                    &    (.00353)         &    (.00253)         &    (.00183)         &    (.00405)         &    (.00304)         &    (.00117)         \\
\addlinespace
4th quartile        &      .00669         &     -.00929\sym{***}&     -.00613\sym{***}&     -.00272         &      .00466         &     -.00386\sym{***}\\
                    &     (.0108)         &     (.0028)         &    (.00214)         &    (.00422)         &    (.00319)         &    (.00132)         \\
\addlinespace
Constant            &        .223\sym{***}&        .238\sym{***}&        .245\sym{***}&        .264\sym{***}&         .22\sym{***}&         .24\sym{***}\\
                    &     (.0142)         &       (.01)         &    (.00796)         &     (.0123)         &     (.0136)         &    (.00474)         \\
\midrule
Observations        &       91212         &       97955         &      170788         &       56419         &       80708         &      498842         \\
\bottomrule
\multicolumn{7}{l}{\footnotesize Standard errors in parentheses}\\
\multicolumn{7}{l}{\footnotesize \sym{*} \(p<0.1\), \sym{**} \(p<0.05\), \sym{***} \(p<0.01\)}\\
\end{tabular}
}
}
        {{
\def\sym#1{\ifmmode^{#1}\else\(^{#1}\)\fi}
\begin{tabular}{l*{6}{c}}
\toprule
                    &\multicolumn{1}{c}{Adivasi}&\multicolumn{1}{c}{Dalit}&\multicolumn{1}{c}{OBC}&\multicolumn{1}{c}{Forward}&\multicolumn{1}{c}{Muslim}&\multicolumn{1}{c}{\shortstack{All five\\social groups}}\\
\midrule
\textbf{Predictors of gestation ≥3 months}&                     &                     &                     &                     &                     &                     \\
\addlinespace
\hspace*{1em}Not using modern contraception&       0.080\sym{***}&       0.072\sym{***}&       0.062\sym{***}&       0.055\sym{***}&       0.091\sym{***}&       0.071\sym{***}\\
                    &    (0.0015)         &    (0.0013)         &    (0.0009)         &    (0.0012)         &    (0.0018)         &    (0.0006)         \\
\addlinespace
\hspace*{1em}Less than primary education&      -0.011\sym{***}&      -0.001         &       0.004\sym{***}&      -0.000         &      -0.005\sym{*}  &      -0.001         \\
                    &    (0.0017)         &    (0.0014)         &    (0.0010)         &    (0.0013)         &    (0.0022)         &    (0.0006)         \\
\addlinespace
\hspace*{1em}Rural resident      &       0.010\sym{***}&       0.006\sym{***}&       0.004\sym{***}&       0.008\sym{***}&       0.009\sym{***}&       0.006\sym{***}\\
                    &    (0.0027)         &    (0.0017)         &    (0.0012)         &    (0.0015)         &    (0.0024)         &    (0.0008)         \\
\addlinespace
\hspace*{1em}Does not have boy child&       0.016\sym{***}&       0.022\sym{***}&       0.026\sym{***}&       0.023\sym{***}&       0.018\sym{***}&       0.022\sym{***}\\
                    &    (0.0025)         &    (0.0024)         &    (0.0017)         &    (0.0021)         &    (0.0033)         &    (0.0010)         \\
\addlinespace
\textbf{Age categories} \\ (15–19 omitted)&                     &                     &                     &                     &                     &                     \\
\addlinespace
\hspace*{1em}20–24               &       0.001         &       0.002         &       0.010         &      -0.003         &       0.024\sym{*}  &       0.008\sym{*}  \\
                    &    (0.0084)         &    (0.0078)         &    (0.0065)         &    (0.0111)         &    (0.0096)         &    (0.0037)         \\
\addlinespace
\hspace*{1em}25–29               &      -0.024\sym{**} &      -0.034\sym{***}&      -0.034\sym{***}&      -0.035\sym{**} &      -0.002         &      -0.028\sym{***}\\
                    &    (0.0085)         &    (0.0079)         &    (0.0064)         &    (0.0108)         &    (0.0097)         &    (0.0037)         \\
\addlinespace
\hspace*{1em}30–49               &      -0.076\sym{***}&      -0.075\sym{***}&      -0.075\sym{***}&      -0.077\sym{***}&      -0.058\sym{***}&      -0.074\sym{***}\\
                    &    (0.0083)         &    (0.0078)         &    (0.0064)         &    (0.0107)         &    (0.0097)         &    (0.0037)         \\
\addlinespace
\textbf{Parity \& time since last live birth categories} \\ (No prior births omitted)&                     &                     &                     &                     &                     &                     \\
\addlinespace
\hspace*{1em}1 birth, below 2y spacing&      -0.159\sym{***}&      -0.117\sym{***}&      -0.127\sym{***}&      -0.140\sym{***}&      -0.190\sym{***}&      -0.140\sym{***}\\
                    &    (0.0063)         &    (0.0062)         &    (0.0047)         &    (0.0068)         &    (0.0080)         &    (0.0028)         \\
\addlinespace
\hspace*{1em}1 birth, 2–3y spacing&      -0.060\sym{***}&      -0.032\sym{***}&      -0.048\sym{***}&      -0.088\sym{***}&      -0.098\sym{***}&      -0.060\sym{***}\\
                    &    (0.0083)         &    (0.0082)         &    (0.0061)         &    (0.0080)         &    (0.0106)         &    (0.0036)         \\
\addlinespace
\hspace*{1em}1 birth, above 3y spacing&      -0.078\sym{***}&      -0.080\sym{***}&      -0.080\sym{***}&      -0.100\sym{***}&      -0.095\sym{***}&      -0.085\sym{***}\\
                    &    (0.0059)         &    (0.0056)         &    (0.0042)         &    (0.0057)         &    (0.0082)         &    (0.0025)         \\
\addlinespace
\hspace*{1em}2 births, below 2y spacing&      -0.153\sym{***}&      -0.131\sym{***}&      -0.141\sym{***}&      -0.152\sym{***}&      -0.197\sym{***}&      -0.149\sym{***}\\
                    &    (0.0063)         &    (0.0061)         &    (0.0044)         &    (0.0063)         &    (0.0081)         &    (0.0027)         \\
\addlinespace
\hspace*{1em}2 births, 2–3y spacing&      -0.093\sym{***}&      -0.090\sym{***}&      -0.090\sym{***}&      -0.124\sym{***}&      -0.122\sym{***}&      -0.099\sym{***}\\
                    &    (0.0078)         &    (0.0068)         &    (0.0052)         &    (0.0069)         &    (0.0099)         &    (0.0031)         \\
\addlinespace
\hspace*{1em}2 births, above 3y spacing&      -0.101\sym{***}&      -0.099\sym{***}&      -0.100\sym{***}&      -0.115\sym{***}&      -0.125\sym{***}&      -0.105\sym{***}\\
                    &    (0.0056)         &    (0.0052)         &    (0.0039)         &    (0.0055)         &    (0.0078)         &    (0.0023)         \\
\addlinespace
\hspace*{1em}3+ births, below 2y spacing&      -0.137\sym{***}&      -0.123\sym{***}&      -0.124\sym{***}&      -0.135\sym{***}&      -0.158\sym{***}&      -0.130\sym{***}\\
                    &    (0.0061)         &    (0.0059)         &    (0.0046)         &    (0.0073)         &    (0.0084)         &    (0.0027)         \\
\addlinespace
\hspace*{1em}3+ births, 2–3y spacing&      -0.066\sym{***}&      -0.066\sym{***}&      -0.074\sym{***}&      -0.097\sym{***}&      -0.109\sym{***}&      -0.075\sym{***}\\
                    &    (0.0075)         &    (0.0073)         &    (0.0057)         &    (0.0087)         &    (0.0097)         &    (0.0033)         \\
\addlinespace
\hspace*{1em}3+ births, above 3y spacing&      -0.103\sym{***}&      -0.098\sym{***}&      -0.098\sym{***}&      -0.113\sym{***}&      -0.144\sym{***}&      -0.105\sym{***}\\
                    &    (0.0055)         &    (0.0052)         &    (0.0039)         &    (0.0055)         &    (0.0076)         &    (0.0023)         \\
\addlinespace
\textbf{Wealth quartiles} \\ (1st i.e. bottom quartile omitted)&                     &                     &                     &                     &                     &                     \\
\addlinespace
\hspace*{1em}2nd quartile        &       0.001         &      -0.004\sym{*}  &      -0.003\sym{*}  &      -0.002         &      -0.004         &      -0.003\sym{**} \\
                    &    (0.0018)         &    (0.0018)         &    (0.0015)         &    (0.0026)         &    (0.0029)         &    (0.0009)         \\
\addlinespace
\hspace*{1em}3rd quartile        &       0.002         &      -0.004\sym{*}  &      -0.004\sym{*}  &      -0.000         &      -0.001         &      -0.003\sym{**} \\
                    &    (0.0023)         &    (0.0019)         &    (0.0015)         &    (0.0025)         &    (0.0031)         &    (0.0009)         \\
\addlinespace
\hspace*{1em}4th quartile        &      -0.001         &      -0.007\sym{**} &      -0.001         &       0.002         &      -0.002         &      -0.003\sym{***}\\
                    &    (0.0033)         &    (0.0022)         &    (0.0017)         &    (0.0026)         &    (0.0034)         &    (0.0010)         \\
\addlinespace
\hspace*{1em}Constant            &       0.156\sym{***}&       0.153\sym{***}&       0.152\sym{***}&       0.167\sym{***}&       0.167\sym{***}&       0.157\sym{***}\\
                    &    (0.0093)         &    (0.0082)         &    (0.0067)         &    (0.0112)         &    (0.0105)         &    (0.0039)         \\
\midrule
\textbf{N}          &      91,205         &      97,944         &     170,783         &      80,703         &      56,415         &     497,050         \\
\bottomrule
\end{tabular}
}
}

        \begin{tablenotes}
            \item Standard errors are in parentheses and are clustered at the primary sampling unit level. Social groups are defined based on caste and religion as follows: Adivasi (ST), Dalit (SC), Muslim (non‐SC/ST), OBC (Hindu/Sikh OBC), and Forward caste (Hindus not SC/ST/OBC). Sample includes married women who are either currently pregnant or not using modern contraception, and are in one of these social groups. 
            \item * $p<0.5$, ** $p<0.01$
        \end{tablenotes}
    \end{threeparttable}
\end{table}

\begin{landscape}
    


\begin{table}[H]
    \centering
    \small % shrink text

    \begin{threeparttable}
        
    
    \caption{: Sample sizes of 3+ month pregnant women in each subgroup and the \% of pregnant women dropped to estimate pre-pregnancy underweight}
    \label{tab:sampledropped}
    \renewcommand{\arraystretch}{1.5}
    \resizebox{\textwidth}{!}{%
        \begin{tabular}{l*{12}{>{\centering\arraybackslash}p{1.2cm}}}
\toprule
& \multicolumn{2}{c}{\shortstack{All five \\\\ social groups}} & \multicolumn{2}{c}{Adivasi} & \multicolumn{2}{c}{Dalit} & \multicolumn{2}{c}{OBC} & \multicolumn{2}{c}{Forward} & \multicolumn{2}{c}{Muslim} \\\\
\cmidrule(lr){2-3} \cmidrule(lr){4-5} \cmidrule(lr){6-7} \cmidrule(lr){8-9} \cmidrule(lr){10-11} \cmidrule(lr){12-13}
Predictor Group & N & \% dropped & N & \% dropped & N & \% dropped & N & \% dropped & N & \% dropped & N & \% dropped \\\\
\midrule
\textbf{Parity and time since last live birth categories}&&&&&&&&&&&&\\
\hspace*{2em}No births&9,841&0.04&1,855&0.05&1,968&0.10&3,251&0.03&1,414&0.00&1,353&0.00\\
\hspace*{2em}1 birth, \textless{}2y spacing&1,650&1.45&244&1.64&424&0.94&617&1.62&207&1.93&158&1.27\\
\hspace*{2em}1 birth, 2–3y spacing&2,097&1.38&407&1.47&470&1.91&741&0.81&243&0.41&236&2.97\\
\hspace*{2em}1 birth, \textgreater{}3y spacing&3,999&0.13&879&0.11&666&0.45&1,217&0.08&642&0.00&595&0.00\\
\hspace*{2em}2 births, \textless{}2y spacing&738&0.95&145&3.45&201&0.00&238&0.00&64&1.56&90&1.11\\
\hspace*{2em}2 births, 2–3y spacing&1,027&1.07&219&0.00&214&0.93&354&0.56&83&2.41&157&3.18\\
\hspace*{2em}2 births, \textgreater{}3y spacing&2,019&0.35&567&0.71&328&0.30&433&0.23&208&0.00&483&0.21\\
\hspace*{2em}3+ births, \textless{}2y spacing&573&0.87&138&1.45&126&1.59&143&0.00&39&2.56&127&0.00\\
\hspace*{2em}3+ births, 2–3y spacing&944&1.69&277&1.44&213&1.41&233&0.00&56&12.50&165&1.21\\
\hspace*{2em}3+ births, \textgreater{}3y spacing&1,193&0.75&404&0.00&218&1.38&258&0.78&67&1.49&246&1.22\\
 &&&&&&&&&&&&\\
\textbf{Wealth Categories}&&&&&&&&&&&&\\
\hspace*{2em}1st quartile&7,032&0.68&2,423&0.58&1,564&1.02&1,749&0.34&366&1.09&930&0.86\\
\hspace*{2em}2nd quartile&6,242&0.34&1,458&0.21&1,343&0.30&1,983&0.25&599&1.00&859&0.35\\
\hspace*{2em}3rd quartile&5,767&0.50&888&0.68&1,134&0.53&2,005&0.35&812&0.37&928&0.75\\
\hspace*{2em}4th quartile&5,040&0.38&366&1.09&787&0.38&1,748&0.29&1,246&0.32&893&0.34\\
 &&&&&&&&&&&&\\
\textbf{All predictor groups}&24,081&0.49&5,135&0.53&4,828&0.60&7,485&0.31&3,023&0.56&3,610&0.58\\
\bottomrule
\end{tabular}

    }
    
    \begin{tablenotes}
    \footnotesize
    \item \parbox{\textwidth}{
        Pregnant women are dropped in reweighting if they are categorized into bins based on the reweighting covariates with no nonpregnant women. 
    }
    \end{tablenotes}

    \end{threeparttable}
\end{table}

\end{landscape}



% \begin{landscape}

% \subsection{Outcomes and percent dropped for different reweighting specifications}
\newpage

\begin{landscape}
    


\begin{table}[H]
    \centering

    \begin{threeparttable}
            
    \caption{: Prepregnancy underweight and diagnostics from different reweighting specifications}
    \label{specifications}
    \renewcommand{\arraystretch}{1.5}
    \scriptsize
    \setlength{\tabcolsep}{5pt}
    \resizebox{\textwidth}{!}{%
        \begin{tabular}{l*{5}{>{\centering\arraybackslash}p{1.5cm}}}
\toprule
 & \multicolumn{5}{c}{Proportion of prepregnancy underweight} \\\\
\cmidrule(lr){2-6}
Variables used in reweighting relative to original specification & Adivasi & Dalit & OBC & Forward & Muslim \\\\
\midrule
Original specification (4 category age bins)$^a$&0.27&0.23&0.20&0.16&0.14\\
Replace age bins to 5 category$^b$&0.27&0.23&0.21&0.16&0.14\\
Replace age bins to 5 category$^b$, remove education&0.27&0.23&0.21&0.15&0.15\\
Replace age bins to 5 category$^b$, remove rural&0.27&0.23&0.21&0.15&0.14\\
Replace age bins to 5 category$^b$, remove rural and education&0.27&0.23&0.20&0.15&0.14\\
Replace age bins to 6 category$^c$, remove rural and education&0.27&0.23&0.21&0.15&0.15\\
Replace age bins to 7 category$^d$, remove rural and education&0.27&0.23&0.21&0.15&0.15\\
\bottomrule
\end{tabular}

    }
    
    \scriptsize

    
    \begin{tablenotes}
        \item  \textit{Notes}: columns titled Adivasi, Dalit, OBC, Forward and Muslim depict  prepregnancy underweight in those subgroups using the given reweighting specification.
        \item $^a$ Original reweighting specification includes all categories seen in Table \ref{tab:sumstats} and used in the main results
        \item $^a$ 4 category age bin (15–19, 20–24, 25–29, 30–49), 10 parity \& time since last birth categories, 4 wealth quartiles, and indicators for not using modern contraception, less than primary education, rural resident, and does not have boy child.
        \item $^b$ 5 category age bin has 15-19, 20-24, 25-29, 30-34, and 35-49.
        \item $^c$ 6 category age bin has 15-19, 20-22, 23-25, 26-29, 30-34, and 35-49.
        \item $^d$ 7 category age bin has 15-17, 18-20, 21-23, 24-26, 27-30, and 31-34 and 35-39.
        \item $^e$ decomposition subgroups refers to wealth quartile or parity and birth spacing subgroups within each social group (70 subgroups total)

        
    \end{tablenotes}

    
    
    \end{threeparttable}
\end{table}

% \end{landscape}

\end{landscape}

% \end{landscape}
\newpage
\begin{table}[H]
    \centering

    \begin{threeparttable}
            
    \caption{: Fertility rates by social group}
    \label{tab:tfr}
    \renewcommand{\arraystretch}{1.5}
    \scriptsize
    \setlength{\tabcolsep}{5pt}
    \resizebox{\textwidth}{!}{%
        % \begin{tabular}{l>{\centering\arraybackslash}p{1.4cm}}
\toprule
Social group & TFR \\\\
\midrule
Adivasi&2.147991\\
Dalit&2.133749\\
OBC&2.018258\\
Forward&1.677633\\
Muslim&2.420128\\
All five social groups&2.047685\\
\bottomrule
\end{tabular}

        \begin{tabular}{l*{6}{>{\centering\arraybackslash}p{1.9cm}}}
\toprule
ASFR & Adivasi & Dalit & OBC & Forward & Muslim & All five social groups \\\\
\midrule
15-19&0.14&0.14&0.14&0.13&0.15&0.14\\
20-24&0.22&0.23&0.23&0.19&0.23&0.22\\
25-29&0.16&0.16&0.16&0.14&0.17&0.16\\
30-34&0.10&0.09&0.07&0.06&0.10&0.08\\
35-39&0.04&0.03&0.02&0.02&0.04&0.03\\
40-44&0.01&0.01&0.00&0.00&0.01&0.01\\
45-49&0.00&0.00&0.00&0.00&0.00&0.00\\
TFR&3.41&3.30&3.10&2.68&3.47&3.16\\
95\% CI for TFR&[3.35, 3.46]&[3.25, 3.36]&[3.06, 3.14]&[2.63, 2.74]&[3.40, 3.54]&[3.14, 3.18]\\
\bottomrule
\end{tabular}

    }
    
    \small
    \begin{tablenotes}
        \item  \textit{Notes:} As in the DHS reports, fertility rates are based on births in the three years before the survey. Social groups are defined based on self-reported constitutional categories and religion as follows: Adivasi (ST), Dalit (SC), Muslim (non‐SC/ST), OBC (Hindu/Sikh OBC), and Forward caste (Hindus not SC/ST/OBC).
    \end{tablenotes}

    \end{threeparttable}
\end{table}

\newpage

\section{S.\arabic{section} Use of AI} \label{sec:aiprompts}
In Section \ref{sec:AIuse}, we describe how we used AI in the production of this research.  We write:

\begin{quote}
We used OpenAI ChatGPT (GPT-5.2) to assist with Stata programming, table and figure formatting, and editing. Use cases included (i) requests to rewrite existing Stata code to improve runtime, (ii) requests to generate code in the Stata packages listtex and esttab to achieve desired table layouts, (iii) assistance with navigating Stata documentation to achieve desired figure layouts, and (iv) copyediting to reduce word count and improve clarity.
\end{quote}

Below, we provide sample prompts for each use case.

Prompt for use case (i):
\emph{This is my dofile that does prepregnancy reweighting, it creates a tag for every possible combination of the reweighting variables and then collapses to get the count of pregnant and nonpregnant women in each bin to flag the bins that have only nonpregnant women. It works but it takes too long in my 1000 iteration bootstrap, please suggest how I can make it more efficient.} 

Prompt for use case (ii):
\emph{I want to generate the LaTeX code for a table in Stata that shows the percent of the sample that is in each subgroup category, keeping rows as all the levels of my reweighting variables and columns as each of the 5 social groups, pregnant and nonpregnant. This is my dofile where I'm currently doing estadd scalar for dummy models, but it doesn't allow me to create header spacing in the rownames for similar categories - is there a way I can just edit my table as a .dta file in the Stata console directly and then export to LaTeX without esttab?}

Prompt for use case (iii): 
\emph{I have a collapsed results dataset where rows are subgroups and there is a column for mean, lower level, and upper level of the confidence interval. How do I plot these to where subgroups are categories, non-numerical, on the x axis each with one scatter point with confidence interval bars?}

Prompt for use case (iv):
\emph{Here is my work-in-progress relevance statement, please suggest places where I should cut words to make it 100 words.}

\newpage

\section{S.\arabic{section} Additional maternal nutrition indicators by social group} \label{sec:weightgain}

\subsection{S.\arabic{section}.\arabic{subsection} Description of additional maternal nutrition indicators} \label{sec:weightgainm1}

\subsection{S.\arabic{section}.\arabic{subsection} Weight gain method 1} \label{sec:weightgainm1}

The first method of estimating weight gain starts by estimating monthly weight gain of pregnant women at 3+ months gestational duration as $\hat{\beta}$ from an OLS regression of the form
\begin{equation}
    weight_i = \alpha + \beta month_i + \tau C_i + \epsilon_i
\end{equation}

where $weight_i$ is the weight in kilograms of pregnant woman $i$. $month_i$ is month of gestation, calculated as described in the methods section, and $C_i$ are controls, including women's current ages (single year), years of education, number of living children, urban or rural residence, wealth quintile, as well as state-month fixed effects.

Assuming women gain 10 times more weight in their second and third trimesters than in their first, we estimate weight gain for a full-term pregnancy as 
\begin{equation}
    \hat{weight gain} = 1.1\times6\times\hat{\beta}
\end{equation}

In table A1, the 95\% confidence intervals for weight gain method 1 are calculated as $1.1\times6\times\hat{\beta} \pm 1.96\times se_{\hat{\beta}}$.

\subsection{S.\arabic{section}.\arabic{subsection} Weight gain method 2} \label{sec:weightgainm2}


The second method of estimating weight gain takes the difference of the average weight of pregnant women at gestational duration of 9 months or more and the average prepregnancy weight using the nonparametric reweighting method. 


\begin{equation}
    \hat{weight gain} = (\bar{W}_9-\bar{W}_0)+(0.5)\hat{\beta}
\end{equation}
where $\bar{W}_9$ is the average weight of pregnant women at gestational duration of 9 months or more, and $\bar{W}_0$ is the average prepregnancy weight using the nonparametric reweighting method. To account for additional weight gain between the time of survey and the end of pregnancy for women of gestational duration of 9 months or more, we assume these women are at the midpoint of their last month of pregnancy and add the accordingly scaled $\hat{\beta}$ from the regression in weight gain method 1.

This offers a benchmark to method 1, which may be inaccurate if the assumption is false that women gain 10 times more weight in their second and third trimesters than in their first, or if gestational durations are incorrectly reported.





\newpage
\begin{landscape}
\begin{table}[!htbp]
\centering
\begin{adjustbox}{center, max width=0.95\linewidth} % <-- reduced slightly from 1.0 to 0.95
\begin{threeparttable}
    \small
    \caption{: Means and 95\% CIs for additional maternal nutrition indicators by social group}
    \label{tab:addnut}

    \setlength{\tabcolsep}{2pt}
    \renewcommand{\arraystretch}{1.35}

    \begin{tabular}{l*{6}{>{\centering\arraybackslash}p{3.5cm}}}
\toprule
\Social Group & BMI & Weight (lbs) & Overweight^a & Obesity^b & Weight gain \newline method 1 (lbs)^c & Weight gain \newline method 2 (lbs)^d \\\\
\midrule
Adivasi&20.6 [20.5, 20.8]&103.8 [102.8, 104.9]&0.13 [0.11, 0.14]&0.05 [0.05, 0.06]&15.0 [12.4, 17.5]&14.8 [11.9, 18.8]\\
Dalit&21.2 [21.1, 21.3]&106.5 [106.1, 107.1]&0.18 [0.17, 0.19]&0.07 [0.07, 0.08]&15.5 [13.3, 17.7]&15.4 [13.5, 17.1]\\
OBC&21.5 [21.4, 21.5]&109.5 [109.0, 110.0]&0.20 [0.19, 0.21]&0.08 [0.07, 0.08]&16.6 [14.7, 18.5]&15.5 [13.9, 17.1]\\
Forward&22.2 [22.1, 22.4]&114.7 [114.0, 115.5]&0.26 [0.25, 0.27]&0.11 [0.10, 0.12]&18.3 [14.8, 21.7]&14.3 [10.8, 17.2]\\
Muslim&22.1 [22.0, 22.1]&112.8 [112.3, 113.4]&0.25 [0.24, 0.26]&0.11 [0.10, 0.11]&11.4 [7.9, 14.9]&12.7 [8.9, 15.8]\\
All five social groups&21.5 [21.5, 21.6]&109.5 [109.4, 110.0]&0.20 [0.20, 0.21]&0.08 [0.08, 0.09]&15.6 [14.4, 16.8]&14.7 [13.5, 15.8]\\
\bottomrule
\end{tabular}


    \begin{tablenotes}[para,flushleft]
        \item Estimates are generated using a nonparametric reweighting function applied to nonpregnant women in the original sample (married, not using modern contraception, and belonging to one of the social groups of focus). Confidence intervals are computed from resampling clusters from the original sample, stratifying by type of place of residence (urban/rural) and state.
        \item $^{a}$ Overweight defined as BMI above 25.
        \item $^{b}$ Obesity defined as BMI above 30.
        \item $^{c}$ Weight gain method 1 and $^{d}$ weight gain method 2 explained above.
        \item Social groups are defined based on self-reported constitutional categories and religion as follows: Adivasi (ST), Dalit (SC), Muslim (non‐SC/ST), OBC (Hindu/Sikh OBC), and Forward caste (Hindus not SC/ST/OBC).
    \end{tablenotes}

\end{threeparttable}
\end{adjustbox}
\end{table}
\end{landscape}



\section{S.\arabic{section} Alt text for figures}

\textbf{Figure 1:}  An interval plot of the prevalence of prepregnancy underweight within Adivasi, Dalit, OBC, forward caste Hindu, and Muslim women in India, and for all groups together. (167 characters)

\textbf{Figure 2:} Six-panel figure. Top three show stacked bar charts of pregnant women by parity, birth spacing, and wealth. Bottom three show interval plots of prepregnancy underweight across the same categories. (196 characters)

\textbf{Figure A1:} A bar chart with confidence intervals of the proportion of women who report eating last in their household within each Adivasi, Dalit, OBC, forward caste Hindu, and Muslim social groups in India.


 
\end{document}
